% -----------------------------------------------------------------------------
% chapters/chap-graph.tex —— 图论案例：图的表示、度与最短路
%
% 本文件演示在章节中按需使用较重的宏包：画图需要 TikZ。
% 模板的类文件出于“精简依赖”的考虑不加载 TikZ，需要时请在主文件的导言区
% 写 \usepackage{tikz}；这样只有真正需要画图的文档才会付出编译代价。
% -----------------------------------------------------------------------------
\chapter{图论案例：图的表示、度与最短路}

\section{图与记号}

\begin{definition}[图]
设 $V$ 是一个非空有限集合，$E$ 是 $V$ 的若干二元子集构成的集合。
称有序对 $(V, E)$ 为一个\textbf{图}，记作 $G = (V, E)$，其中 $V$ 是图的
\textbf{顶点集}，$E$ 是图的\textbf{边集}；$E$ 中的每个二元子集称为一条边。
\end{definition}

\begin{remark}
若边 $\{u, v\}$ 的两个端点可以相同，则称 $u$ 与 $v$ 之间存在\textbf{自环}；
若 $E$ 中不含自环且任意两个顶点至多只相邻于一条边，则称 $G$ 为\textbf{简单图}。
本文讨论的图均为简单无向图。
\end{remark}

\section{度、路径与连通性}

\begin{definition}[度]
顶点 $v$ 的\textbf{度} $d(v)$ 是与 $v$ 相邻的顶点数；记 $\deg(G) = \sum_{v \in V} d(v)$。
\end{definition}

\begin{theorem}[握手定理]
任意有限图 $G$ 满足 $\sum_{v \in V} d(v) = 2|E|$。
\end{theorem}

\begin{proof}
对每条边计数：它连接两个顶点，对两个顶点的度数各贡献 $1$，
故对所有顶点的度数求和时每条边被计入两次。
\end{proof}

\begin{definition}[路径与连通]
顶点序列 $v_0, v_1, \dots, v_k$ 称为从 $v_0$ 到 $v_k$ 的一条\textbf{长度}为 $k$ 的
\textbf{路径}，若 $v_{i-1}$ 与 $v_i$ 相邻且各顶点互不相同。
若 $v_0 = v_k$，则称其为一个\textbf{圈}。若 $G$ 中任意两个顶点之间都存在路径，
则称 $G$ 是\textbf{连通}的。
\end{definition}

\section{图的一个例子}

考虑图 $G = (V, E)$，其中 $V = \{v_1, v_2, v_3, v_4, v_5\}$，
边集为 $E = \{\{v_1,v_2\}, \{v_2,v_3\}, \{v_1,v_4\}, \{v_4,v_5\}, \{v_5,v_3\}, \{v_2,v_5\}\}$，
各条边的权值如图 \ref{fig:graph-weighted} 所示。

\begin{figure}[htbp]
  \centering
  % 本图需要 TikZ，类文件不加载，请见主文件导言区的 \usepackage{tikz}
  \begin{tikzpicture}[
    scale=1.3,
    vertex/.style={circle,draw=structurecolor,fill=structurecolor!12,
                   inner sep=2pt,minimum size=6mm},
    weight/.style={fill=white,inner sep=1pt,font=\scriptsize}]
    \node[vertex] (v1) at (0,0) {$v_1$};
    \node[vertex] (v2) at (2.4,0) {$v_2$};
    \node[vertex] (v3) at (4.8,0) {$v_3$};
    \node[vertex] (v4) at (1.1,1.9) {$v_4$};
    \node[vertex] (v5) at (3.7,1.9) {$v_5$};
    \draw (v1) -- node[weight,above] {$1$} (v2);
    \draw (v2) -- node[weight,above] {$3$} (v3);
    \draw (v1) -- node[weight,left] {$2$} (v4);
    \draw (v4) -- node[weight,above] {$4$} (v5);
    \draw (v5) -- node[weight,right] {$1$} (v3);
    \draw (v2) -- node[weight,below right] {$5$} (v5);
  \end{tikzpicture}
  \caption{一个带权无向图}
  \label{fig:graph-weighted}
\end{figure}

\begin{exercise}
写出图 \ref{fig:graph-weighted} 中各顶点的度，并判断该图是否连通。
\end{exercise}

\section{图的表示}

计算机中存储图有两种常见方式：一是\textbf{邻接矩阵}，二是\textbf{邻接表}。
邻接矩阵用 $a_{ij} = 1$ 表示 $v_i$ 与 $v_j$ 相邻，用 $0$ 表示不相邻；
对于带权图，则把 $1$ 换成边权，不相邻的位置记为 $\infty$。

\begin{table}[htbp]
  \centering
  \caption{图 \ref{fig:graph-weighted} 的邻接矩阵}
  \label{tab:adjacency}
  \begin{tabular}{c|ccccc}
    \toprule
    & $v_1$ & $v_2$ & $v_3$ & $v_4$ & $v_5$ \\
    \midrule
    $v_1$ & 0 & 1 & 0 & 1 & 0 \\
    $v_2$ & 1 & 0 & 1 & 0 & 1 \\
    $v_3$ & 0 & 1 & 0 & 0 & 1 \\
    $v_4$ & 1 & 0 & 0 & 0 & 1 \\
    $v_5$ & 0 & 1 & 1 & 1 & 0 \\
    \bottomrule
  \end{tabular}
\end{table}

邻接矩阵需要 $O(|V|^2)$ 的存储空间，稠密图上运算方便；邻接表只需 $O(|V| + |E|)$ 的空间，
在稀疏图上更高效。

\section{最短路与遍历算法}

\begin{definition}[最短路]
在带权图中，路径 $P$ 的\textbf{长度}是其上各条边权之和。
顶点 $s$ 到 $t$ 的\textbf{最短路}是所有 $s$ 到 $t$ 的路径中长度最小者。
\end{definition}

\begin{custom}{算法：广度优先搜索}
从起点 $s$ 出发，逐层扩展：先访问所有与 $s$ 相邻的顶点，再访问与它们相邻的顶点，
如此进行，直到所有可达顶点都被访问。记录每个顶点首次被访问时所经过的边，
即得到一棵\textbf{广度优先搜索树}。
\end{custom}

\begin{note}
若所有边的权值相同（或者都取 $1$），广度优先搜索得到的就是最短路；
若边权各不相同，则需要 Dijkstra 算法等加权算法。
\end{note}

\begin{exercise}
在图 \ref{fig:graph-weighted} 上执行一次广度优先搜索，写出搜索树，
并求 $v_1$ 到 $v_3$ 的最短路长度。
\end{exercise}
